% =====================================================================
% Drop-in section: A Composition Theorem for Pipeline Deletion
% ---------------------------------------------------------------------
% INTEGRATION NOTES
% 1. Add these declarations to the PREAMBLE of main.tex (you already
% have \newtheorem{proposition}{Proposition}; do NOT redeclare it):
%
% \newtheorem{theorem}{Theorem}
% \newtheorem{lemma}{Lemma}
% \newtheorem{corollary}{Corollary}
% \theoremstyle{definition}
% \newtheorem{definition}{Definition}
% \newtheorem{assumption}{Assumption}
%
% 2. \input this file where the formal contribution should appear,
% e.g. just after Section 4 (Deletion-Risk Metric) or in the
% supplement. It uses your \method{} and \bench{} macros.
% 3. Requires amsmath, amssymb/amsfonts, amsthm (already loaded).
% =====================================================================

\section{A Composition Theorem for Pipeline Deletion}
\label{sec:composition}

The experiments show repeatedly that deleting a record from each
component in isolation does not delete it from the deployed system: an
index can be purged while a summary survives, or an adapter can be
unlearned while a derivative is still retrieved. This section
makes that observation precise. We define behavioral deletion against a
counterfactual reference system, prove that per-component deletion
guarantees \emph{compose} up to an explicit derivative-coupling term,
and prove that this coupling term is \emph{necessary}: even perfect,
independently certified component-level deletion does not imply
pipeline-level deletion. The theorem provides the formal justification
for treating deletion verification as a provenance-closed, multi-channel
behavioral audit, which is exactly what \method{} operationalizes.

\subsection{Setup and Counterfactual Reference}

Fix a deletion target $z$. A foundation-model pipeline produces, for a
query $q$, a distribution over answers. We write the pipeline's answer
law as $P_{S}(\cdot\mid q)$ for a system $S$. A system is characterized
by a parameter state $\theta$ (base weights and adapters) and an
external knowledge state $K$ (vector index, lexical index, reranker
state, caches, summaries, and any -derivative store). We
consider queries drawn from a declared probe class $\mathcal{B}$ and
measure deviation between answer laws by total variation,
$\mathrm{TV}(\mu,\nu)=\sup_{A}\lvert\mu(A)-\nu(A)\rvert$.

We compare three systems built from the same architecture:
\begin{itemize}
\item $S_0$, the pre-deletion system $(\theta_0,K_0)$;
\item $S^{-}$, the deployed post-deletion system $(\theta^{-},K^{-})$
obtained by applying a deletion operation to $S_0$;
\item $S^{\star}$, the \emph{counterfactual reference} $(\theta^{\star},
K^{\star})$, i.e.\ the pipeline that would have been built had $z$
\emph{never been ingested}. $S^{\star}$ is the gold standard inherited
from exact retraining and certified removal, lifted here from a single
model to the whole pipeline.
\end{itemize}

\begin{definition}[Behavioral residual influence]
\label{def:residual}
The behavioral residual influence of $z$ in the post-deletion system,
relative to the probe class $\mathcal{B}$, is
\[
 \rho_{\mathcal{B}}(S^{-})
 \;=\;
 \sup_{q\in\mathcal{B}}
 \mathrm{TV}\!\big(P_{S^{-}}(\cdot\mid q),\,P_{S^{\star}}(\cdot\mid q)\big).
\]
We say $S^{-}$ achieves \emph{$\varepsilon$-behavioral deletion} on
$\mathcal{B}$ if $\rho_{\mathcal{B}}(S^{-})\le\varepsilon$.
\end{definition}

Definition~\ref{def:residual} is the property an auditor actually cares
about: it asks whether the deployed system is behaviorally
indistinguishable, over the probe class, from a system that never saw
$z$. It is strictly stronger than checking that artifacts labeled with
$z$ are gone, a distinction we now name.

\begin{definition}[Syntactic vs.\ behavioral deletion]
\label{def:syntactic}
A deletion operation is \emph{syntactic} if it removes exactly the
artifacts explicitly tagged with the provenance identifier of $z$. It is
\emph{behavioral} if the resulting system attains
$\varepsilon$-behavioral deletion for small $\varepsilon$. Syntactic
deletion is observable from metadata; behavioral deletion is a property
of the answer law and is what Definition~\ref{def:residual} measures.
\end{definition}

\subsection{Pipeline Decomposition}

We model the canonical retrieval-plus-generation pipeline as two stages
with the parameter state entering generation as a side input. Given $q$
and knowledge state $K$, a retrieval stage produces a (random) context
$X$ with law $\mathcal{R}(\cdot\mid q,K)$. A generation stage produces
the answer with law $\mathcal{G}(\cdot\mid q,X,\theta)$. The pipeline
answer law is the mixture
\begin{equation}
\label{eq:pipeline}
 P_{S}(\cdot\mid q)
 \;=\;
 \mathbb{E}_{X\sim\mathcal{R}(\cdot\mid q,K)}
 \big[\,\mathcal{G}(\cdot\mid q,X,\theta)\,\big].
\end{equation}
Deeper pipelines (reranking, caching, summarization) are obtained by
further composing stages; the serial generalization is
Theorem~\ref{thm:general}.

We assume the deletion operation is, or can be analyzed relative to, a
\emph{provenance-closed} idealization. Let $\mathcal{D}(z)$ denote the
transitive closure of $z$ in the provenance graph: every chunk,
embedding, summary, cache entry, instruction/QA example, and
fine-tuning record causally derived from $z$ or from a member of
$\mathcal{D}(z)$. Let $S^{\diamond}=(\theta^{\diamond},K^{\diamond})$ be
the system obtained from $S^{-}$ by \emph{additionally} removing all of
$\mathcal{D}(z)$ from the knowledge state and excluding it from the data
that produced the parameters. $S^{\diamond}$ is the system $S^{-}$
``should have been'' had deletion followed lineage rather than labels.

\subsection{Assumptions}

\begin{assumption}[Parametric behavioral deletion]
\label{a:theta}
Uniformly over the probe class and over contexts,
\[
 \sup_{q\in\mathcal{B},\,x}
 \mathrm{TV}\!\big(\mathcal{G}(\cdot\mid q,x,\theta^{\diamond}),\,
 \mathcal{G}(\cdot\mid q,x,\theta^{\star})\big)\le\varepsilon_{\theta}.
\]
\end{assumption}

\begin{assumption}[Knowledge-state deletion]
\label{a:k}
For a retrieval discrepancy $d_{\mathcal{R}}$,
\[
 \sup_{q\in\mathcal{B}}
 d_{\mathcal{R}}\!\big(\mathcal{R}(\cdot\mid q,K^{\diamond}),\,
 \mathcal{R}(\cdot\mid q,K^{\star})\big)\le\varepsilon_{K}.
\]
\end{assumption}

\begin{assumption}[Bounded context sensitivity]
\label{a:lip}
The generation stage is $L$-Lipschitz from the retrieval discrepancy to
answer total variation: for any context laws $\alpha,\beta$ and any
fixed $(q,\theta)$,
\[
 \mathrm{TV}\!\Big(\mathbb{E}_{X\sim\alpha}\mathcal{G}(\cdot\mid q,X,\theta),\,
 \mathbb{E}_{X\sim\beta}\mathcal{G}(\cdot\mid q,X,\theta)\Big)
 \le L\, d_{\mathcal{R}}(\alpha,\beta).
\]
\end{assumption}

Assumptions~\ref{a:theta}--\ref{a:k} say the parametric and knowledge
components are individually deletion-certified \emph{against the
counterfactual references on the provenance-closed problem}.
Assumption~\ref{a:lip} bounds how strongly a change in what is retrieved
can change what is answered. When $d_{\mathcal{R}}$ is itself total
variation over context laws, $L\le 1$ by the data-processing inequality
(mixing is a channel), and \eqref{eq:pipeline} gives no amplification;
$L>1$ captures the practical case where $\varepsilon_{K}$ is reported in
a weaker retrieval metric (e.g.\ retrieved-identifier overlap) than the
answer divergence it induces.

\subsection{Composition}

\begin{lemma}[Hybrid bound]
\label{lem:hybrid}
For any $q$,
\[
 \mathrm{TV}\!\big(P_{S^{\diamond}}(\cdot\mid q),P_{S^{\star}}(\cdot\mid q)\big)
 \le \varepsilon_{\theta} + L\,\varepsilon_{K}.
\]
\end{lemma}

\begin{proof}
Insert the hybrid system $(\theta^{\star},K^{\diamond})$ and apply the
triangle inequality:
\[
\begin{aligned}
\mathrm{TV}\!\big(P_{(\theta^{\diamond},K^{\diamond})},P_{(\theta^{\star},K^{\star})}\big)
&\le
\mathrm{TV}\!\big(P_{(\theta^{\diamond},K^{\diamond})},P_{(\theta^{\star},K^{\diamond})}\big)\\
&\quad+
\mathrm{TV}\!\big(P_{(\theta^{\star},K^{\diamond})},P_{(\theta^{\star},K^{\star})}\big).
\end{aligned}
\]
For the first term the knowledge state is fixed at $K^{\diamond}$, so by
\eqref{eq:pipeline} and convexity of total variation under mixtures,
\[
\begin{aligned}
&\mathrm{TV}\!\big(P_{(\theta^{\diamond},K^{\diamond})},
P_{(\theta^{\star},K^{\diamond})}\big)\\
&\quad\le
\mathbb{E}_{X}\,
\mathrm{TV}\!\big(\mathcal{G}(\cdot\mid q,X,\theta^{\diamond}),
\mathcal{G}(\cdot\mid q,X,\theta^{\star})\big)
\le \varepsilon_{\theta},
\end{aligned}
\]
using Assumption~\ref{a:theta}. For the second term the parameters are
fixed at $\theta^{\star}$ and only the retrieval law differs, so by
Assumption~\ref{a:lip} and then Assumption~\ref{a:k},
\[
\begin{aligned}
&\mathrm{TV}\!\big(P_{(\theta^{\star},K^{\diamond})},
P_{(\theta^{\star},K^{\star})}\big)\\
&\quad\le L\,d_{\mathcal{R}}\!\big(
\mathcal{R}(\cdot\mid q,K^{\diamond}),
\mathcal{R}(\cdot\mid q,K^{\star})\big)
\le L\,\varepsilon_{K}.
\end{aligned}
\]
Summing gives the claim.
\end{proof}

We now isolate the part of the deviation that per-component guarantees
\emph{cannot} reach: the orphan derivatives present in $S^{-}$ but absent
from the provenance-closed $S^{\diamond}$.

\begin{definition}[Derivative-coupling term]
\label{def:gamma}
\[
 \gamma_{z}
 \;=\;
 \sup_{q\in\mathcal{B}}
 \mathrm{TV}\!\big(P_{S^{-}}(\cdot\mid q),\,P_{S^{\diamond}}(\cdot\mid q)\big).
\]
We say the deletion operation achieves \emph{provenance closure} if
$S^{-}=S^{\diamond}$, equivalently if every artifact in
$\mathcal{D}(z)$ was removed.
\end{definition}

\begin{theorem}[Pipeline deletion composition]
\label{thm:main}
Under Assumptions~\ref{a:theta}--\ref{a:lip},
\[
 \boxed{\;\rho_{\mathcal{B}}(S^{-})
 \;\le\;
 \varepsilon_{\theta}\;+\;L\,\varepsilon_{K}\;+\;\gamma_{z}.\;}
\]
Moreover $\gamma_{z}=0$ if and only if the deletion operation achieves
provenance closure, in which case
$\rho_{\mathcal{B}}(S^{-})\le\varepsilon_{\theta}+L\,\varepsilon_{K}$.
\end{theorem}

\begin{proof}
By the triangle inequality, for each $q$,
\[
\mathrm{TV}\!\big(P_{S^{-}},P_{S^{\star}}\big)
\le
\mathrm{TV}\!\big(P_{S^{-}},P_{S^{\diamond}}\big)
+\mathrm{TV}\!\big(P_{S^{\diamond}},P_{S^{\star}}\big).
\]
Taking $\sup_{q\in\mathcal{B}}$, the first supremum is $\gamma_{z}$ by
Definition~\ref{def:gamma} and the second is bounded by
$\varepsilon_{\theta}+L\,\varepsilon_{K}$ by Lemma~\ref{lem:hybrid}.
If provenance closure holds then $S^{-}=S^{\diamond}$, so
$\gamma_{z}=0$; conversely $\gamma_{z}=0$ forces
$P_{S^{-}}=P_{S^{\diamond}}$ on $\mathcal{B}$, i.e.\ the retained
derivatives exert no behavioral influence over the probe class, which is
provenance closure in behavioral form.
\end{proof}

The next result shows the coupling term is not an artifact of loose
analysis: component-wise deletion, \emph{even when perfect}, is
insufficient.

\begin{proposition}[Necessity of provenance closure]
\label{prop:necessity}
For every $\delta\in(0,1)$ there is a pipeline and a deletion operation
with $\varepsilon_{\theta}=\varepsilon_{K}=0$---each component
individually certified against its counterfactual reference and free of
any artifact \emph{labeled} $z$---yet
$\rho_{\mathcal{B}}(S^{-})\ge 1-\delta$. Hence no bound of the form
$\rho_{\mathcal{B}}(S^{-})\le f(\varepsilon_{\theta},\varepsilon_{K})$
can hold without the coupling term $\gamma_{z}$.
\end{proposition}

\begin{proof}
Let the base model never be trained on $z$, so
$\theta^{-}=\theta^{\diamond}=\theta^{\star}$ and
$\varepsilon_{\theta}=0$. During ingestion a summarization component
writes a derivative $s(z)\in\mathcal{D}(z)$---a paraphrased summary of
$z$'s protected fact---into the index under a \emph{fresh} artifact
identifier that is not tagged with $z$. Syntactic deletion removes the
raw chunk of $z$ but, having no label linking $s(z)$ to $z$, leaves
$s(z)$ in $K^{-}$. Restrict the comparison to direct artifacts of $z$:
relative to those, the retrieval laws of $K^{-}$ and $K^{\star}$
coincide, so a label-based audit reports $\varepsilon_{K}=0$.

Now take a probe $q\in\mathcal{B}$ asking for the protected fact. Under
$S^{-}$ the retriever returns $s(z)$ with probability one and the
generator answers the fact; under $S^{\star}$ neither $z$ nor $s(z)$ was
ever ingested, so the system returns ``insufficient information.'' These
answer laws are almost disjoint, giving
$\mathrm{TV}(P_{S^{-}}(\cdot\mid q),P_{S^{\star}}(\cdot\mid q))\ge1-\delta$
for any target $\delta$ by making the two responses distinguishable with
probability $\ge1-\delta$. Therefore $\rho_{\mathcal{B}}(S^{-})\ge1-\delta$
while $\varepsilon_{\theta}=\varepsilon_{K}=0$; the entire deviation is
carried by $\gamma_{z}$, the orphan derivative $s(z)$.
\end{proof}

Proposition~\ref{prop:necessity} is the formal counterpart of the
\emph{ derivative retained} and \emph{shadow-copy} failures in
our experiments, where each component looks clean by its own check yet
the pipeline still answers the deleted fact. It explains why those rows
carry near-maximal risk even though no artifact labeled with the target
remains.

\subsection{Serial Pipelines}

\begin{theorem}[Multi-stage composition]
\label{thm:general}
Consider a serial pipeline $C_{m}\circ\cdots\circ C_{1}$ in which stage
$C_{i}$ has deletion error $\varepsilon_{i}$ against its counterfactual
reference (on the provenance-closed problem) and every downstream stage
$C_{k}$, $k>i$, is $L_{k}$-Lipschitz in the sense of
Assumption~\ref{a:lip}. Then
\[
 \rho_{\mathcal{B}}(S^{-})
 \;\le\;
 \sum_{i=1}^{m}\Big(\textstyle\prod_{k=i+1}^{m}L_{k}\Big)\varepsilon_{i}
 \;+\;\gamma_{z}.
\]
\end{theorem}

\begin{proof}
Form the $m$ hybrids that switch stages from the reference to the
deployed version one at a time, holding all others fixed. Telescoping
the triangle inequality, the deviation introduced when stage $C_{i}$ is
switched is at most $\varepsilon_{i}$, and it is transported to the
output through the downstream stages, each contracting (or amplifying)
it by its Lipschitz factor, contributing
$(\prod_{k>i}L_{k})\varepsilon_{i}$. Summation over $i$ bounds the
provenance-closed deviation $\mathrm{TV}(P_{S^{\diamond}},P_{S^{\star}})$;
adding $\gamma_{z}$ as in Theorem~\ref{thm:main} gives the result.
\end{proof}

\subsection{Consequences for Auditing}

\begin{corollary}[Audit sufficiency]
\label{cor:audit}
To certify $\varepsilon$-behavioral deletion it suffices to (i) establish
provenance closure, driving $\gamma_{z}\to 0$, and (ii) bound each
component error so that the propagated sum in
Theorem~\ref{thm:general} is at most $\varepsilon$. Conversely, by
Proposition~\ref{prop:necessity}, an audit that omits either ingredient
cannot certify deletion.
\end{corollary}

Corollary~\ref{cor:audit} maps onto \method{} term by term. Building the
provenance graph and expanding deletion targets to $\mathcal{D}(z)$ is
the operational route to $\gamma_{z}\to 0$; the residual-dependence
channels estimate the component and behavioral errors that survive
closure (direct, paraphrase, and extraction probes estimate the
answer-side parametric and derivative error; retrieval-dependence probes
estimate the retrieval-mediated term $L\varepsilon_{K}$; watermark
probes detect labeled-artifact survival; the retrieval-off counterfactual
isolates which path dominates). The Hoeffding upper confidence bound of
Section~4 is then the finite-sample, high-probability estimate of
$\rho_{\mathcal{B}}(S^{-})$ that Theorem~\ref{thm:main} bounds in
expectation: a pass certifies that the estimated propagated error plus
coupling term lies below tolerance, under the declared probe class.

\paragraph{Relation to differential-privacy composition.}
Theorem~\ref{thm:general} is structurally a composition-across-data-flow
result---deletion budgets add (sensitivity-weighted) along the
pipeline---but it differs from differential-privacy composition in an
essential way: it carries the non-additive term $\gamma_{z}$. Standard
privacy composition has no analogue of $\gamma_{z}$ because it does not
model derived artifacts that re-enter the pipeline. The coupling term is
precisely the contribution of derivative re-ingestion, and it is why
deletion in compositional foundation-model systems requires lineage
tracking rather than per-component accounting alone.

\paragraph{Assumptions and scope.}
The guarantee is conditional on the declared probe class $\mathcal{B}$
(it bounds residual influence over $\mathcal{B}$, not over all
conceivable prompts), on access to the counterfactual references for
defining $\varepsilon_{\theta},\varepsilon_{K}$ (in practice estimated,
e.g.\ via a retrain-without-record reference on a sampled slice), and on
stable system behavior during the audit. These are the same operating
assumptions used by the empirical metric, now made explicit at the level
of the composition they justify.
